\begin{frame}[plain]
  \centering
  \begin{columns}[c]
    \begin{column}{0.46\textwidth}
      \centering
      \includegraphics[height=0.56\textheight]{hochreiter_portrait.jpg}\\[0.3em]
      {\small Sepp Hochreiter \\[0.1em] LSTM 논문 1저자 (1997, 당시 박사과정)}
    \end{column}
    \begin{column}{0.46\textwidth}
      \centering
      \includegraphics[height=0.56\textheight]{schmidhuber_portrait.jpg}\\[0.3em]
      {\small Jürgen Schmidhuber \\[0.1em] 지도교수, LSTM 공동 저자}
    \end{column}
  \end{columns}
  \vskip0.3em
  {\footnotesize Hochreiter \& Schmidhuber, ``Long Short-Term Memory,''
  \textit{Neural Computation} 9(8), 1997 --- 곱 대신 덧셈으로 기억을 전달하자는 해법을 이미 제안}
\end{frame}
